% !TEX encoding = UTF-8
% !TEX program = xelatex
% preamble.tex - LaTeX 模板导言区
% 兼容 Pandoc 转换 Word 与 XeLaTeX 独立编译
%
% 本文件定义了符合学校《撰写格式及要求》的全部样式：
% 页面设置、字体、行距、标题格式、图表标题、公式编号、摘要/致谢/参考文献环境等

\usepackage{fontspec}                 % 字体设置
\setmainfont{Times New Roman}         % 英文字体（学校要求）
% ctexart 文档类已自动加载中文字体，\heiti \songti \fangsong 等命令均可用
% 无需重复定义，避免冲突

\usepackage{amsmath, amssymb}         % 数学公式支持
\usepackage{graphicx}                 % 插图支持
\usepackage[bookmarks=true, colorlinks=false, hidelinks]{hyperref} % 超链接（PDF书签）
\usepackage{geometry}                 % 页面设置
\geometry{
    a4paper,
    top=2.2cm, bottom=2.2cm,
    left=2.5cm, right=2.5cm,
    headheight=1.2cm, footskip=1.5cm   % 页眉页脚高度，符合学校要求
}
\usepackage{fancyhdr}                 % 页眉页脚
\pagestyle{fancy}
\fancyhf{}
\fancyhead[C]{\zihao{-5}\heiti XeTeX正文内容模板---pdf产物仅供预览---需通过模板转换} % 页眉居中，小5号黑体
\fancyfoot[C]{\thepage}               % 页码居中
\renewcommand{\headrulewidth}{0pt}    % 去掉页眉横线

\usepackage{setspace}                 % 行距控制
\onehalfspacing                        % 全文 1.5 倍行距（学校规范）
\usepackage{makecell}                 % 支持单元格内手动换行，仅保障pdf输出，需要人工设置ref.docx
\usepackage{titlesec}                 % 标题格式自定义
\usepackage{titletoc}                 % 目录格式控制
\usepackage{caption}                  % 图表标题格式
\usepackage{booktabs}                 % 三线表支持
\usepackage{array}                    % 表格列格式增强
\usepackage{float}                    % 浮动体 [H] 参数
\usepackage{enumitem}                 % 列表环境控制
\usepackage[perpage]{footmisc}        % 脚注每页重置

% --- 表格增强设置 ---
% 调整表格行高，使内容更舒展
\renewcommand{\arraystretch}{1.3}
% --------------------

% ---------- 图片格式 ----------
% XeLaTeX 编译用 png（\includegraphics 不支持 svg），
% Pandoc 转 Word 时可改为 svg（Word 原生支持 svg 嵌入）
% 注意：不能用 \renewcommand{\includegraphics} 做内联替换，
%       Pandoc 展开不了 TeX 宏，会导致 Word 中所有图片丢失。
%       仅有的两个 .svg 引用（02.introduce.tex:14, 05.systematic_design.tex:20）
%       编译 XeLaTeX 前手动改为 .png 即可。
\newcommand{\imageformat}{png}

% ---------- 段落缩进 ----------
\setlength{\parindent}{2em}           % 每段首行缩进2字符（学校规范）

% ---------- 目录 ----------
\setcounter{tocdepth}{3}              % 目录显示到三级标题

% 目录条目格式：宋体小4号（中文），Times New Roman 小4号（数字/英文），不倾斜
\titlecontents{section}
    [0em]
    {\songti\zihao{-4}}
    {\thecontentslabel\ }
    {}
    {\dotfill\ \thecontentspage}
    []
\titlecontents{subsection}
    [2em]
    {\songti\zihao{-4}}
    {\thecontentslabel\ }
    {}
    {\dotfill\ \thecontentspage}
    []
\titlecontents{subsubsection}
    [4em]
    {\songti\zihao{-4}}
    {\thecontentslabel\ }
    {}
    {\dotfill\ \thecontentspage}
    []

% ---------- 标题格式（严格遵循学校规范） ----------
% 一级标题：小2号黑体，加粗，居中，新页开始，序号后空1字符
\titleformat{\section}
    {\newpage\centering\heiti\zihao{-2}\bfseries}
    {\thesection}{1em}{}
\titlespacing*{\section}{0pt}{0pt}{0pt}

% 二级标题：小3号黑体，靠左顶格
\titleformat{\subsection}
    {\heiti\zihao{-3}\bfseries}
    {\thesubsection}{1em}{}
\titlespacing*{\subsection}{0pt}{0pt}{0pt}

% 三级标题：4号黑体，靠左顶格
\titleformat{\subsubsection}
    {\heiti\zihao{4}\bfseries}
    {\thesubsubsection}{1em}{}
\titlespacing*{\subsubsection}{0pt}{0pt}{0pt}

% 四级及以下标题（如需要）可用 \paragraph 等，学校未严格要求，此处省略

% ---------- 图表标题格式（学校规范） ----------
% 要求：5号宋体加粗，图序与图题之间空1字符，段前段后1行（约12pt）
\captionsetup{
    font={bf,small},          % small ≈ 5号(10.5pt)，bf 加粗
    labelfont={bf,small},
    textfont={bf,small},
    labelsep=space,           % 序号与标题间空一个字符
    skip=12pt,                % 与上下文间距（段前段后1行）
    position=bottom           % 图标题在下，表标题在上（由环境自动处理）
}

% ---------- 图表与公式编号：按章排序（如”图1-1”、”表2-3”、”(3-2)”） ----------
\usepackage{chngcntr}
\counterwithin{figure}{section}
\counterwithin{table}{section}
\counterwithin{equation}{section}
\renewcommand{\thefigure}{\thesection--\arabic{figure}}
\renewcommand{\thetable}{\thesection--\arabic{table}}
\renewcommand{\theequation}{\thesection--\arabic{equation}}

% ---------- 目录标题中文化 ----------
\renewcommand{\contentsname}{目 \qquad 录}

% ---------- 中文摘要环境 ----------
\newenvironment{cnabstract}
    {\newpage
     \phantomsection
     \addcontentsline{toc}{section}{\texorpdfstring{摘 \quad 要}{摘 要}}
     \begin{center}
        \heiti\zihao{-2}\bfseries 摘 \quad 要
     \end{center}
     \vspace{12pt}
     \zihao{-4}\songti
    }{\par\vspace{12pt}}

% ---------- 英文摘要环境 ----------
\newenvironment{enabstract}
    {\newpage
     \phantomsection
     \addcontentsline{toc}{section}{Abstract}  % 纯英文无需 \texorpdfstring
     \begin{center}
        \bfseries\zihao{-2} Abstract
     \end{center}
     \vspace{12pt}
     \zihao{-4}\rmfamily
    }{\par\vspace{12pt}}

% ---------- 关键词命令 ----------
\newcommand{\cnkeywords}[1]{%
    \par\noindent{\heiti\zihao{-4}\bfseries 关键词：}\songti\zihao{-4} #1
}
\newcommand{\enkeywords}[1]{%
    \par\noindent{\bfseries\zihao{-4} Key words:} \zihao{-4} #1
}

% ---------- 致谢环境 ----------
\newenvironment{thesisthanks}
    {\newpage
     \phantomsection
     \addcontentsline{toc}{section}{\texorpdfstring{致 \quad 谢}{致 谢}}
     \begin{center}
        \heiti\zihao{-2}\bfseries 致 \quad 谢
     \end{center}
     \vspace{12pt}
     \zihao{-4}\songti
    }{\par}

% ---------- 参考文献环境 ----------
% 使用手动的 thebibliography，字体为仿宋小四
\newenvironment{thesisbibl}
    {\newpage
     \phantomsection
     \addcontentsline{toc}{section}{参考文献}
     \begin{center}
        \heiti\zihao{-2}\bfseries 参考文献
     \end{center}
     \vspace{12pt}
     \zihao{-4}\fangsong
     \setlength{\parindent}{0pt}      % 参考文献条目悬挂缩进（需在 bib 条目中自行控制）
     \begin{thebibliography}{ }} %换行确保格式一致性
    {\end{thebibliography}\par}

% ---------- 附录环境 ----------
\newenvironment{thesisappendix}
    {\newpage
     \phantomsection
     \addcontentsline{toc}{section}{附 \quad 录}
     \begin{center}
        \heiti\zihao{-2}\bfseries 附 \quad 录
     \end{center}
     \vspace{12pt}
     \zihao{-4}\songti
    }{\par}

% ---------- 防止浮动体跨章节 ----------
\usepackage[section]{placeins}
